\chapter{Introducción}

\section{Objetivos generales}

El planteamiento inicial de este Trabajo de Fin de Grado partió de una
propuesta del profesor Eduardo Manuel Sánchez Vila, tutor del trabajo,
dirigida a su grupo de investigación. Todo nació a partir de un TFM que había llevado el propio Eduardo.
Un TFM que sentó las bases del entorno de simulación, tanto en cuestión de código como de idea. Una idea original que explotaríamos varios niveles por encima para crear un sistema nuevo, vivo y complejo.
La propuesta original perseguía
sentar las bases de un entorno común sobre el que estudiar paradigmas de
toma de decisiones humanas. A partir de esa idea, el alcance, la
arquitectura y la división del proyecto en dos fases se fueron concretando
de forma iterativa entre Eduardo, mi compañero y yo, en sucesivas reuniones de
diseño que reorientaron el enfoque inicial, redefiniendo requisitos y enriqueciendo el producto.

El proyecto resultante es un laboratorio virtual para la simulación y el
análisis de paradigmas de toma de decisiones humanas mediante agentes
inteligentes, dividido en dos fases complementarias. La primera,
desarrollada por Pablo Pazos Parada, implementa un pipeline que transforma
descripciones en lenguaje natural en modelos de decisión ejecutables. Esta
memoria documenta la segunda, desarrollada por Juan Freire Álvarez:
la infraestructura que recoge esos modelos, los ejecuta en un entorno
común, observa su comportamiento y sintetiza los hallazgos en informes
comparativos.

El objetivo general de esta fase es construir una infraestructura que
pueda ejecutar modelos de decisión heterogéneos y producir, en cada
simulación, un conjunto de artefactos verificables:

\begin{itemize}
\item una especificación del entorno de simulación adaptada al problema
descrito;
\item un registro estructurado del comportamiento de los modelos durante
la ejecución;
\item un análisis de los patrones observados, contrastado con postulados
previos del Knowledge Backbone;
\item un informe en PDF con los resultados y la evidencia recogida;
\item un conjunto de memorias de simulación persistidas para alimentar
ejecuciones futuras.
\end{itemize}

Estos artefactos se obtienen mediante un pipeline de agentes
especializados coordinados por un Orchestrator, y se exponen al
usuario a través de una interfaz web. La justificación de cada uno de
estos componentes se desarrolla en los capítulos correspondientes.

La integración con la primera fase se realiza por \textit{duck
typing} sobre el contrato \texttt{decide}, \texttt{update} y
\texttt{get\_state}, sin acoplamiento de clases entre ambos sistemas.

\section{Descripción del sistema}

El sistema es un laboratorio virtual para ejecutar y comparar paradigmas
de toma de decisiones humanas. Se accede a través de una interfaz web --- y,
de forma secundaria, mediante una línea de comandos para uso interno y de
depuración --- y se opera desde un chat con un agente coordinador, el
Orchestrator. En cada sesión, el Orchestrator dirige un pipeline compuesto
por cuatro agentes especializados (Architect, Tracker, Analyst y Reporter):
crea un entorno de simulación adaptado al problema, ejecuta los modelos de
decisión, observa y analiza el comportamiento resultante, y produce un
informe en PDF.

Los modelos de decisión que se ejecutan no son parte de esta fase, sino
que se cargan dinámicamente desde la primera fase del proyecto, que es la
encargada de generarlos. La integración entre ambas fases se realiza a
través de un contrato mínimo --- los métodos \texttt{decide},
\texttt{update} y \texttt{get\_state} --- aplicado por \textit{duck
typing} en el momento de la carga.

El acceso al Knowledge Backbone compartido entre fases es
bidireccional: la segunda fase consume desde él tanto los propios
modelos disponibles para simular como las memorias previas que la
primera fase consolida sobre cada paradigma (sus formulaciones, su
procedencia y las predicciones de comportamiento asociadas), y devuelve
las observaciones registradas por el Tracker para que las ejecuciones
posteriores puedan apoyarse en ellas.

\section{Información adicional de interés}

La arquitectura del sistema parte de varias decisiones técnicas que se
resumen a continuación y se desarrollan con más detalle en el capítulo de
diseño.

\paragraph{Multi-agente coordinado por un Orchestrator centralizado.} El
sistema podría haberse implementado como un único agente, pero dividir
las responsabilidades en cuatro agentes especializados hace cada salida
intermedia inspeccionable y reduce la superficie de error. El contraste
con la alternativa monolítica se desarrolla en el capítulo de estado
del arte.

\paragraph{Acceso al conocimiento mediado por el Orchestrator.} El
acceso al Knowledge Backbone se centraliza: el Analyst dispone de la
herramienta de recuperación como \textit{tool call} y la invoca cuando
necesita ampliar el contexto, mientras que el Architect y el Reporter
no consultan el backbone por su cuenta: reciben un contexto
pre-recuperado por el propio Orchestrator antes de cada invocación. El
detalle del esquema y de sus consecuencias se desarrolla en el capítulo
de diseño.

\paragraph{Backend FastAPI con un único canal WebSocket.} El cliente web
debe reaccionar a la actividad de los agentes a medida que se produce,
sin recurrir a \textit{polling}. Un canal WebSocket bidireccional,
complementado por unos pocos \textit{endpoints} REST, simplifica el
código del backend y la sincronización del cliente.

\paragraph{Knowledge Backbone híbrido.} La capa de persistencia combina
cuatro componentes con responsabilidades diferenciadas --- metadatos
relacionales, artefactos binarios, recuperación vectorial y léxica, y
grafo de paradigmas ---. Los capítulos de requisitos y diseño explican
qué datos residen en cada uno y por qué se separan de ese modo.

\paragraph{Integración con la primera fase por \textit{duck typing}.}
En lugar de compartir un paquete con clases base entre las dos fases,
el contrato se reduce a tres métodos. La verificación se realiza en
tiempo de ejecución mediante \texttt{runtime\_checkable}. Esto evita
acoplar la evolución de ambas fases y hace que un cambio en una no
obligue a reorganizar la otra.

\section{Relación de la documentación}

La memoria sigue la estructura de un Trabajo de Fin de Grado de tipo B,
orientado al desarrollo de un sistema software. Los capítulos principales
son:

\begin{itemize}
\item estado del arte;
\item especificación de requisitos;
\item diseño del sistema;
\item desarrollo del trabajo;
\item pruebas;
\item conclusiones y posibles ampliaciones.
\end{itemize}

Los apéndices incluyen el manual técnico, el manual de usuario y la
licencia. La bibliografía recoge tanto referencias sobre agentes
inteligentes y modelos de decisión como trabajos en recuperación aumentada,
grafos de conocimiento y simulación multi-agente.
